The Software Testing Life Cycle (STLC) involves several activities, starting with understanding requirements and planning tests, followed by test design, execution, defect reporting, and release decisions. In most development teams, these activities are handled using different tools, often requiring manual effort to move information from one stage to another. This project explores how AI can help bring these activities together through a central orchestrator supported by a group of specialised agents. The aim is to help testing teams make better use of requirements, user stories, code changes, and test execution results to generate useful quality-related insights, while keeping important decisions, such as reviewing results and approving releases, under human control.

The prototype developed as part of this project is a web-based application built using React and TypeScript for the frontend, FastAPI for the backend, and PostgreSQL for data storage. A separate agent service, developed using Google's Agent Development Kit (ADK), handles AI-related tasks. The backend manages user authentication, stores assistance requests, and communicates with the agent service. At the centre of the system is a root agent that identifies the nature of a request and directs it to one of nine specialised agents. These agents focus on different STLC activities, including test-case generation, test-data generation, test automation suggestions, test execution analysis, self-healing recommendations, change-impact analysis, defect reporting, test-summary preparation, and release-readiness assessment. The application also includes integrations for importing Jira Cloud stories and retrieving code changes from GitHub. It supports optional Microsoft Entra single sign-on and allows different model providers to be configured.

This report covers the motivation behind the project, relevant research on software testing and AI agents, the proposed architecture, implementation details, and security considerations. It also discusses how the system can be evaluated in a practical setting. An important distinction is made between the features currently implemented in the prototype and the capabilities proposed for future development. For example, the test execution analysis agent works with available test evidence but does not execute external test suites. Similarly, self-healing recommendations are provided as suggestions rather than being applied automatically. No completed comparative benchmark or measured performance improvement is claimed in this report.

The prototype brings several testing-related activities into a single application and demonstrates how specialised AI agents can support different stages of the STLC. However, its actual effectiveness, reliability, and suitability for use in real-world projects will need to be assessed through further testing and controlled evaluation. Future evaluation can focus on the quality of generated outputs, traceability, time and effort saved, execution efficiency, and the accuracy of defect analysis, with human review remaining an important part of the process.
